\documentclass{article}
\usepackage[italian]{babel}
\usepackage[T1]{fontenc}
\usepackage[utf8]{inputenc}
\usepackage{amsmath}

\title{Introduzione alla geometria differenziale}
\author{G.~Rossi}
\date{2026}

\begin{document}
\maketitle

\section{Introduzione}

La geometria differenziale studia le propriet\`a delle variet\`a lisce
utilizzando gli strumenti del calcolo infinitesimale.  In questo lavoro
consideriamo le superfici regolari nello spazio euclideo tridimensionale.

Sia $S \subset \mathbb{R}^3$ una superficie regolare e sia $p \in S$ un
punto.  Lo spazio tangente $T_p S$ \`e il sottospazio bidimensionale
di~$\mathbb{R}^3$ che meglio approssima~$S$ nel punto~$p$.

\section{Curvatura gaussiana}

\begin{theorem}[Teorema Egregium di Gauss]
La curvatura gaussiana $K$ di una superficie \`e invariante per isometrie.
Pi\`u precisamente, $K$ dipende solo dalla prima forma fondamentale e dalle
sue derivate.
\end{theorem}

Questo risultato fondamentale mostra che la curvatura \`e una propriet\`a
intrinseca della superficie, misurabile senza fare riferimento allo spazio
ambiente.

\section{Conclusione}

La geometria differenziale fornisce strumenti essenziali per la fisica
matematica e la relativit\`a generale.

\end{document}
